    \section{IDENTIDAD DEL TRIPLE PRODUCTO DE JACOBI}

    Esta sección describe una identidad famosa de la teoría de la función zeta debida a Jacobi.
    El teorema de los números pentagonales de Euler y muchas otras identidades con particiones
    se obtienen como casos especiales de la fórmula de Jacobi.

    \begin{teorema}[Identidad del triple producto de Jacobi]\label{teo:14-6}
        Para $x$ y $z$ complejos con $\abs{x}<1$ y $z\neq0$ tenemos
        \begin{equation}\label{eq:14-13}
            \prod_{n=1}^{\infty}(1-x^{2n})(1+x^{2n-1}z^2)(1+x^{2n-1}z^{-2})=\sum_{m=-\infty}^{\infty}x^{m^2}z^{2m}.
        \end{equation}
    \end{teorema}

    \begin{proof}
        La restricción a $\abs{x}<1$ nos asegura la convergencia absoluta de cada uno de los
        productos $\prod(1-x^{2n})$, $\prod(1+x^{2n-1}z^2)$, $\prod(1+x^{2n-1}z^{-2})$, y de la
        serie de \eqref{eq:14-13}. Además, para cada $x$ con $\abs{x}<1$, la serie y los productos
        convergen uniformemente sobre subconjuntos compactos del $z$-plano siempre que no
        contengan a $z=0$, por lo que cada miembro de \eqref{eq:14-13} es una función analítica de
        $z$ para $z\neq0$. Para $z\neq0$ fijo, la serie y los productos también convergen
        uniformemente para $\abs{x}\leq r<1$, luego representan funciones analíticas de $x$ en el
        disco $\abs{x}<1$.

        Para probar \eqref{eq:14-13} fijamos $x$ y definimos $F(z)$ para $z\neq0$ por la ecuación
        \begin{equation}\label{eq:14-14}
            F(z)=\prod_{n=1}^{\infty}(1+x^{2n-1}z^2)(1+x^{2n-1}z^{-2}).
        \end{equation}
        En primer lugar vemos que $F$ satisface la ecuación funcional
        \begin{equation}\label{eq:14-15}
            xz^2F(xz)=F(z).
        \end{equation}
        De \eqref{eq:14-14} obtenemos
        \begin{align*}
            F(xz)&=\prod_{n=1}^{\infty}(1+x^{2n+1}z^2)(1+x^{2n-3}z^{-2})\\
            &=\prod_{m=2}^{\infty}(1+x^{2m-1}z^2)\prod_{r=0}^{\infty}(1+x^{2r-1}z^{-2}).
        \end{align*}
        Puesto que $xz^2=(1+xz^2)/(1+x^{-1}z^{-2})$, la multiplicación de la última ecuación por
        $xz^2$ da \eqref{eq:14-15}.

        Ahora sea $G(z)$ el miembro de la izquierda de \eqref{eq:14-13}, o sea
        \begin{equation}\label{eq:14-16}
            G(z)=F(z)\prod_{n=1}^{\infty}(1-x^{2n}).
        \end{equation}
        Entonces $G(z)$ también satisface la ecuación funcional \eqref{eq:14-15}. Además, $G(z)$
        es una función par de $z$ que es analítica en todo $z\neq0$, luego admite un desarrollo
        de Laurent de la forma
        \begin{equation}\label{eq:14-17}
            G(z)=\sum_{m=-\infty}^{\infty}a_mz^{2m}
        \end{equation}
        en donde $a_{-m}=a_m$ ya que $G(z)=G(z^{-1})$. (Los coeficientes $a_m$ dependen de $x$.)
        Utilizando la ecuación funcional \eqref{eq:14-15} en \eqref{eq:14-17} obtenemos que los
        coeficientes satisfacen la fórmula recursiva
        \begin{equation*}
            a_m=x^{2m-1}a_{m-1}
        \end{equation*}
        que, al iterarla, nos da
        \begin{equation*}
            a_m=a_0x^{m^2}\qquad\text{para todo }m\geq0
        \end{equation*}
        ya que $1+3+\cdots+(2m-1)=m^2$. Este resultado se verifica también para $m<0$. Luego
        \eqref{eq:14-17} se transforma en
        \begin{equation}\label{eq:14-18}
            G_x(z)=a_0(x)\sum_{m=-\infty}^{\infty}x^{m^2}z^{2m},
        \end{equation}
        en donde hemos escrito $G_x(z)$ en vez de $G(z)$ y $a_0(x)$ en vez de $a_0$ para indicar
        la dependencia de $x$. Observamos que \eqref{eq:14-18} implica que $a_0(x)\to1$ cuando
        $x\to0$. Para terminar la demostración debemos probar que $a_0(x)=1$ para todo $x$.

        Si hacemos $z=e^{\pi i/4}$ en \eqref{eq:14-18} obtenemos
        \begin{equation}\label{eq:14-19}
            \frac{G_x(e^{\pi i/4})}{a_0(x)}=\sum_{m=-\infty}^{\infty}x^{m^2}i^m=\sum_{n=-\infty}^{\infty}(-1)^nx^{(2n)^2}
        \end{equation}
        ya que $i^m=-i^{-m}$ si $m$ es impar. De \eqref{eq:14-18} deducimos que la serie de la
        derecha de \eqref{eq:14-19} es $G_{x^4}(i)/a_0(x^4)$ de donde se sigue la identidad
        \begin{equation}\label{eq:14-20}
            \frac{G_x(e^{\pi i/4})}{a_0(x)}=\frac{G_{x^4}(i)}{a_0(x^4)}.
        \end{equation}
        A continuación probamos que $G_x(e^{\pi i/4})=G_{x^4}(i)$. En efecto, \eqref{eq:14-14} y
        \eqref{eq:14-16} nos dan
        \begin{equation*}
            G_x(e^{\pi i/4})=\prod_{n=1}^{\infty}(1-x^{2n})(1+x^{4n-2}).
        \end{equation*}
        Como cada número par es de la forma $4n$ o $4n-2$ tenemos
        \begin{equation*}
            \prod_{n=1}^{\infty}(1-x^{2n})=\prod_{n=1}^{\infty}(1-x^{4n})(1-x^{4n-2})
        \end{equation*}
        luego
        \begin{align*}
            G_x(e^{\pi i/4})&=\prod_{n=1}^{\infty}(1-x^{4n})(1-x^{4n-2})(1+x^{4n-2})=\prod_{n=1}^{\infty}(1-x^{4n})(1-x^{8n-4})\\
            &=\prod_{n=1}^{\infty}(1-x^{8n})(1-x^{8n-4})(1-x^{8n-4})=G_{x^4}(i).
        \end{align*}
        Por lo tanto \eqref{eq:14-20} implica $a_0(x)=a_0(x^4)$. Si substituimos $x$ por $x^4$,
        $x^{4^2},\dots$, obtenemos
        \begin{equation*}
            a_0(x)=a_0(x^{4^k})\qquad\text{para }k=1,2,\dots
        \end{equation*}
        Pero $x^{4^k}\to0$ cuando $k\to\infty$ y $a_0(x)\to1$ cuando $x\to0$, luego $a_0(x)=1$
        para todo $x$. Esto completa la demostración.
    \end{proof}

    \section{CONSECUENCIAS DE LA IDENTIDAD DE JACOBI}

    Si substituimos $x$ por $x^a$ y $z^2$ por $x^b$ en la identidad de Jacobi obtenemos
    \begin{equation*}
        \prod_{n=1}^{\infty}(1-x^{2na})(1+x^{2na-a+b})(1+x^{2na-a-b})=\sum_{m=-\infty}^{\infty}x^{am^2+bm}.
    \end{equation*}
    Análogamente, si $z^2=-x^b$ obtenemos
    \begin{equation*}
        \prod_{n=1}^{\infty}(1-x^{2na})(1-x^{2na-a+b})(1-x^{2na-a-b})=\sum_{m=-\infty}^{\infty}(-1)^mx^{am^2+bm}.
    \end{equation*}
    Para obtener el teorema de los números pentagonales de Euler hacemos simplemente $a=3/2$ y
    $b=1/2$ en esta última igualdad.

    La fórmula de Jacobi conduce a otra fórmula importante para el cubo del producto de Euler.

    \begin{teorema}\label{teo:14-7}
        Si $\abs{x}<1$ tenemos
        \begin{equation}\label{eq:14-21}
            \prod_{n=1}^{\infty}(1-x^n)^3=\sum_{m=-\infty}^{\infty}(-1)^mmx^{(m^2+m)/2}=\sum_{m=0}^{\infty}(-1)^m(2m+1)x^{(m^2+m)/2}.
        \end{equation}
    \end{teorema}

    \begin{proof}
        Substituimos $z^2$ por $-xz$ en la identidad de Jacobi y obtenemos
        \begin{equation*}
            \prod_{n=1}^{\infty}(1-x^{2n})(1-x^{2n}z)(1-x^{2n-2}z^{-1})=\sum_{m=0}^{\infty}(-1)^mx^{m^2+m}\left(z^m-z^{-m-1}\right).
        \end{equation*}
        Ahora reordenamos los términos de ambos miembros, utilizando las relaciones
        \begin{equation*}
            \prod_{n=1}^{\infty}(1-x^{2n-2}z^{-1})=(1-z^{-1})\prod_{n=1}^{\infty}(1-x^{2n}z^{-1})
        \end{equation*}
        y
        \begin{equation*}
            z^m-z^{-m-1}=(1-z^{-1})(1+z^{-1}+z^{-2}+\cdots+z^{-2m})z^m.
        \end{equation*}
        Anulando el factor $1-z^{-1}$ obtenemos
        \begin{multline*}
            \prod_{n=1}^{\infty}(1-x^{2n})(1-x^{2n}z)(1-x^{2n}z^{-1})\\
            =\sum_{m=0}^{\infty}(-1)^mx^{m^2+m}z^m\left(1+z^{-1}+z^{-2}+\cdots+z^{-2m}\right).
        \end{multline*}
        Si hacemos $z=1$ y substituimos $x$ por $x^{1/2}$ obtenemos \eqref{eq:14-21}.
    \end{proof}

    \section{DIFERENCIACIÓN LOGARÍTMICA DE FUNCIONES GENERADORAS}

    El teorema \ref{teo:14-4} da una fórmula recursiva de $p(n)$. Existen otros tipos de
    fórmulas recursivas de funciones aritméticas que se pueden obtener por medio de la
    diferenciación logarítmica de funciones generadoras. A continuación describimos este método.

    Sea $A$ un conjunto de números positivos dado, y sea $f(n)$ una función aritmética.
    Suponemos que el producto

